%% kapa: exam=1 round=36 year=2025 period=1 form=A subject=civil_law subject_ko=민법 seq=25 number=25
%% source: Q-Net 감정평가사 기출문제 (cst00309) artlSeq=5231525 file=raw/1008/5231525/2025년 제36회 감정평가사 1차시험 1교시 문제지 원본.pdf
%% conversion: kordoc 4.24.51 -> tools/gen1.py (auto); review=proofread-sonnet; flags=none
\begin{question}{25}
점유에 관한 설명으로 옳지 않은 것은? (다툼이 있으면 판례에 따름)
\begin{choices}
  \item 물건을 사실상 지배한다는 것은 물건을 물리적, 현실적으로 지배하는 것만을 의미하는 것은 아니다.
  \item 건물의 소유자가 그 건물을 현실적으로 점거하지 아니한 경우, 그는 건물의 부지가 된 토지를 점유하고 있다고 볼 수 없다.
  \item 공터로 형성되어 공중의 이용에 제공되고 있었던 토지 부분을 공로로 나가는 통로로 사용한 것에 불과하다면 그 사용자가 이를 점유하였다고 볼 수 없다.
  \item 건물 공유자 중 일부만이 그 건물을 점유하고 있더라도 그 건물의 부지는 공유자 전원이 공동으로 점유하고 있는 것이다.
  \item 건물의 유치권자는 그 건물을 점유하는 경우에도 그 건물의 부지 부분을 점유하였다고 볼 수는 없다.
\end{choices}
\end{question}
\answer{2}
